\documentclass[11pt,a4paper]{article}
\usepackage{fontspec,amsmath,amssymb,amsthm,microtype}
\setmainfont{lmroman10-regular.otf}[BoldFont=lmroman10-bold.otf,ItalicFont=lmroman10-italic.otf,BoldItalicFont=lmroman10-bolditalic.otf]
\usepackage{xcolor}
\usepackage[most]{tcolorbox}
\usepackage{mathrsfs}
\providecommand{\lt}{<}
\providecommand{\gt}{>}
\setmonofont{DejaVuSansMono.ttf}[BoldFont=DejaVuSansMono-Bold.ttf,ItalicFont=DejaVuSansMono-Oblique.ttf,BoldItalicFont=DejaVuSansMono-BoldOblique.ttf,Scale=0.86]
\newfontfamily\leanfallback{STIXTwoMath-Regular.otf}[Scale=0.86]
\definecolor{leanrule}{gray}{0.62}
\definecolor{leanback}{gray}{0.962}
\definecolor{leantitle}{gray}{0.22}
\definecolor{leanmuted}{gray}{0.38}
\newtcolorbox{leanbox}[1]{enhanced,breakable,colback=leanback,colframe=leanrule,
  boxrule=0.4pt,arc=0.6mm,left=2.4mm,right=2.4mm,top=1.8mm,bottom=1.8mm,
  before skip=0.8em,after skip=1em,fonttitle=\small\bfseries,coltitle=leantitle,
  colbacktitle=leanback,titlerule=0.3pt,title={#1}}
\newlength{\leancharwidth}
\newenvironment{leancode}{\par\smallskip\begingroup\ttfamily\footnotesize\raggedright
  \settowidth{\leancharwidth}{M}\setlength{\parskip}{0pt}\setlength{\parindent}{0pt}
  \frenchspacing}{\par\endgroup\smallskip}
\newcommand{\leanline}[2]{\hangindent=\dimexpr#1\leancharwidth+4\leancharwidth\relax
  \hangafter=1\noindent\hspace*{#1\leancharwidth}#2\par}
\newcommand{\leansource}[1]{\par{\footnotesize\color{leanmuted}#1}\par}
\newcommand{\leanlabel}[1]{\par\smallskip\noindent{\small\bfseries #1}\par}
\usepackage[margin=25mm]{geometry}
\usepackage[unicode,hidelinks]{hyperref}
\hypersetup{pdfauthor={LangevinLean Project}}
\setlength{\emergencystretch}{3em}
\setlength{\parskip}{0.4em}
\setlength{\parindent}{0pt}
\allowdisplaybreaks[2]
\raggedbottom
\title{\texorpdfstring{The genuine Gibbs score observable}{The genuine Gibbs score observable}}
\author{LangevinLean Project}
\date{}
\begin{document}\maketitle
Draft mathematical exposition; these notes may contain mistakes.


The score \(f=\nabla V\) is constructed as a member of \(L^2(\pi;\mathbb R^d)\). Actual stationary evaluation recovers the previously constructed score process, and the Hilbert-space norm is bounded by \(\beta d\).


\section{\texorpdfstring{Constructing the square-integrable gradient}{Constructing the square-integrable gradient}}


\textbf{Claim.} \(\nabla V\in L^2(\pi)\), where \(\pi\) is the normalized Gibbs law of the actual potential.



\textbf{Proof.} Strong convexity makes the Gibbs partition function finite. The shared Gibbs gradient-moment theorem then proves square integrability of the Lipschitz gradient. Passing to its equivalence class preserves the actual gradient as an almost-everywhere representative.

\begin{leanbox}{Lean: \texttt{actual\_\allowbreak{}gradient\_\allowbreak{}memLp}, \texttt{actualObservableScore}, \texttt{actualObservableScore\_\allowbreak{}coe}}
\leansource{\texttt{actual\_\allowbreak{}gradient\_\allowbreak{}memLp} \quad WassersteinMixing/\allowbreak{}ActualScoreObservable.lean, lines 21–25}
\begin{leancode}
\leanline{0}{theorem\ actual\_gradient\_memLp\ :\ MemLp\ (gradient\ V)\ 2\ (ULA.gibbsMeasure\ V)}
\end{leancode}
\leansource{\texttt{actualObservableScore} \quad WassersteinMixing/\allowbreak{}ActualScoreObservable.lean, lines 28–29}
\begin{leancode}
\leanline{0}{def\ actualObservableScore\ :\ Lp\ (E\ P.dimension)\ 2\ (ULA.gibbsMeasure\ V)\ :=}
\leanline{2}{(actual\_gradient\_memLp\ P\ V\ hV).toLp\ (gradient\ V)}
\end{leancode}
\leansource{\texttt{actualObservableScore\_\allowbreak{}coe} \quad WassersteinMixing/\allowbreak{}ActualScoreObservable.lean, lines 31–33}
\begin{leancode}
\leanline{0}{theorem\ actualObservableScore\_coe\ :}
\leanline{4}{actualObservableScore\ P\ V\ hV\ =ᵐ[ULA.gibbsMeasure\ V]\ gradient\ V}
\end{leancode}
\leanlabel{How the Lean proof follows the argument}
\texttt{actual\_\allowbreak{}gradient\_\allowbreak{}mem\allowbreak{}Lp} specializes the Gibbs moment theorem with the actual potential constants. \texttt{actual\allowbreak{}Observable\allowbreak{}Score} uses the resulting \texttt{to\allowbreak{}Lp} constructor and \texttt{actual\allowbreak{}Observable\allowbreak{}Score\_\allowbreak{}coe} records its precise representative.
\end{leanbox}


Earlier results used here: \cite[Section~3.26.10]{LangevinLean2026ULASC}; \cite[Section~3.26.3]{LangevinLean2026ULASC}; \cite[Section~3.27.7]{LangevinLean2026ULASC}.

\section{\texorpdfstring{Evaluating the actual score without a new identification hypothesis}{Evaluating the actual score without a new identification hypothesis}}


\textbf{Claim.} The stationary evaluation map \(U_n f=f(Y_{nh})\) is an isometry, agrees with global evaluation at \(nh\), and satisfies \(U_nf=\nabla V(Y_{nh})\). Consequently \(\|f\|_{L^2(\pi)}^2\le\beta d\).



\textbf{Proof.} The actual grid state has law \(\pi\), giving the isometry. The proved global/\allowbreak{}grid path identity identifies the two evaluation maps. Pull the gradient representative through this measure-preserving state map to identify its evaluated class with the actual drift class. Apply the actual score moment at one time and use the isometry.

\begin{leanbox}{Lean: \texttt{actualGridEvaluation}, \texttt{actualGridEvaluation\_\allowbreak{}eq\_\allowbreak{}global}, \texttt{actualGridEvaluation\_\allowbreak{}score} and 1 more}
\leansource{\texttt{actualGridEvaluation} \quad WassersteinMixing/\allowbreak{}ActualScoreObservable.lean, lines 36–40}
\begin{leancode}
\leanline{0}{def\ actualGridEvaluation\ (n\ :\ ℕ)\ :}
\leanline{4}{Lp\ (E\ P.dimension)\ 2\ (ULA.gibbsMeasure\ V)\ →ₗᵢ[ℝ]}
\leanline{6}{Lp\ (E\ P.dimension)\ 2\ (actualProbability\ P\ V)\ :=}
\leanline{2}{StationaryObservables.evaluation\ (ULA.gibbsMeasure\ V)\ (actualProbability\ P\ V)}
\leanline{4}{(actualReference\ P\ V\ hV\ n)\ (actualReference\_stationary\ P\ V\ hV\ n)}
\end{leancode}
\leansource{\texttt{actualGridEvaluation\_\allowbreak{}eq\_\allowbreak{}global} \quad WassersteinMixing/\allowbreak{}ActualScoreObservable.lean, lines 42–46}
\begin{leancode}
\leanline{0}{theorem\ actualGridEvaluation\_eq\_global\ (n\ :\ ℕ)\ :}
\leanline{4}{actualGridEvaluation\ P\ V\ hV\ n\ =}
\leanline{6}{CanonicalSemigroup.evaluation\ ({\leanfallback 𝕜}\ :=\ ℝ)\ (F\ :=\ E\ P.dimension)\ P\ V\ hV\ (n\ *\ P.step)}
\end{leancode}
\leansource{\texttt{actualGridEvaluation\_\allowbreak{}score} \quad WassersteinMixing/\allowbreak{}ActualScoreObservable.lean, lines 49–61}
\begin{leancode}
\leanline{0}{theorem\ actualGridEvaluation\_score\ (n\ :\ ℕ)\ :}
\leanline{4}{actualGridEvaluation\ P\ V\ hV\ n\ (actualObservableScore\ P\ V\ hV)\ =\ actualScore\ P\ V\ hV\ n}
\end{leancode}
\leansource{\texttt{actualObservableScore\_\allowbreak{}energy} \quad WassersteinMixing/\allowbreak{}ActualScoreObservable.lean, lines 64–69}
\begin{leancode}
\leanline{0}{theorem\ actualObservableScore\_energy\ :}
\leanline{4}{‖actualObservableScore\ P\ V\ hV‖\ \textasciicircum{}\ 2\ ≤\ P.beta\ *\ P.dimension}
\end{leancode}
\leanlabel{How the Lean proof follows the argument}
\texttt{actual\allowbreak{}Grid\allowbreak{}Evaluation\_\allowbreak{}score} transports the almost-everywhere target representative through the actual stationary map and compares it with \texttt{lp\allowbreak{}Drift\_\allowbreak{}coe}. \texttt{actual\allowbreak{}Observable\allowbreak{}Score\_\allowbreak{}energy} transfers the actual probability-space moment through the evaluation isometry.
\end{leanbox}


Earlier results used here: \cite[Section~3.13.3]{LangevinLean2026ULASC}; \cite[Section~3.26.3]{LangevinLean2026ULASC}.

\begin{thebibliography}{ULA-LSI}
\bibitem[ULA-SC]{LangevinLean2026ULASC}
LangevinLean Project. \emph{Unadjusted Langevin under Strong Convexity}. LL-01, 2026. Draft report, content revision d34ead06987c. \url{https://odakyildiz.github.io/LangevinLean/docs/pdf/UnadjustedLangevin/FullProof.pdf}.
\end{thebibliography}

\end{document}
